\chapter*{Kết luận}
\addcontentsline{toc}{chapter}{Kết luận}

ConceptTester chọn đầu vào từ đặc tả và mã của những PUT đã vượt test cơ sở. EP/BVA thử đại diện của lớp và các giá trị sát biên. DT chọn các tổ hợp điều kiện có thể làm thay đổi hành động. ControlFlow tìm dòng hoặc cạnh còn thiếu, dùng quan hệ def--use khi có để giải thích giá trị cần đưa tới đích. Một PUT chỉ được tính khi cùng input cho sai khác ở cả hai lượt chấm.

Một chuỗi có tab cho thấy tác dụng của việc chọn đích trong mã. Trên cùng một bài, nó làm lộ tám PUT mà chỉ ControlFlow Qwen phát hiện. Trên 1.454 PUT đủ điều kiện, hợp ba nhánh cho 121 phát hiện với Qwen và 132 với Gemini. Gộp qua cả hai mô hình, pipeline tìm thêm 44 PUT ngoài Naive LLM. Baseline vẫn có 15 PUT mà pipeline bỏ sót.

Base test loại 169 trong 1.634 PUT ban đầu. Khi chấm thêm trên 162 bài đủ mười cài đặt, test có hướng dẫn từ hai mô hình phát hiện thêm 149 PUT từng đạt. Trên cùng bộ mã Qwen, pass@1 giảm từ 89,75\% xuống 80,56\%. Nếu bỏ mười phát hiện phụ thuộc vào cách hiểu Boolean trong một bài, mức giảm còn 8,58 điểm phần trăm. Những sai khác này cho thấy bộ test cơ sở còn chấp nhận các cài đặt cần kiểm tra thêm.

Mục~\ref{sec:ablation} nêu cách đo riêng từng nhánh hoặc từng kỹ thuật trong EP/BVA với cùng ngân sách. Độ phủ catalogue còn cần bằng chứng rằng input thuộc đúng lớp, biên hoặc quy tắc đã khai báo. Ở các bài có nhiều đáp án hợp lệ, đặc tả vẫn là căn cứ để xác định một sai khác có phải lỗi của PUT hay không.
